
\subsubsection{6.1 Key Findings}

\textbf{Finding 1: SSL shortcut encoding is severe and geometrically stable.}
The 73.36 pp WGA gap between DINO SSL (8.57\%) and supervised training (81.93\%) on Waterbirds is the primary empirical result of this work. More informative than the gap's magnitude is its stability: DINO's WGA does not improve as training progresses, despite converging average accuracy. The best validation WGA (24.06\% at epoch 17) is not maintained --- the final validation WGA (12.03\% at epoch 100) is lower. This rules out the explanation that longer training would close the gap --- the geometry of the SSL loss landscape fixes the shortcut at an early training stage and maintains it. This finding motivates the core proposition: if shortcut encoding is geometrically stable, the intervention must be geometric and must occur during pretraining.

\textbf{Finding 2: The average accuracy / WGA divergence is the diagnostic signature.}
DINO achieves 62.10\% average accuracy and 8.57\% WGA --- a within-model gap of 53.5 pp. This divergence suggests that DINO's self-distillation objective, combined with pretrained ViT-S/8 weights, is effective at encoding background texture as a highly reliable predictor, while the average accuracy metric conceals this systematic failure.

\textbf{Finding 3: The geometric hypothesis remains unverified but well-motivated.}
The sharpness anisotropy measurement (central to H-E1's mechanistic test) is pending the 200-epoch SimCLR/Waterbirds run. We can confirm the existence and severity of the SSL shortcut problem (RQ1 answered). We cannot yet confirm whether that problem is geometrically characterized by directional sharpness anisotropy (RQ2 pending), nor whether SAM during pretraining addresses it (RQ3 pending).

\subsubsection{6.2 Limitations}

\textbf{Limitation 0: Pretrained-weight confound (DINO).}
DINO uses publicly released ImageNet-pretrained ViT-S/8 weights rather than training from scratch on Waterbirds. The 8.57\% WGA may partly reflect background biases encoded during ImageNet pretraining, not only SSL-SGD dynamics on the spurious correlation benchmark. Consequently, the 73.36 pp WGA gap should not be interpreted as arising entirely from SSL training on Waterbirds --- it is an upper bound on the SSL-SGD shortcut effect. The full 9-combination evaluation includes SimCLR and MoCo-v2 trained from scratch on Waterbirds, which will isolate the Waterbirds-specific SSL shortcut effect.

\textbf{Limitation 1: Anisotropy measurement pending.}
The core mechanistic claim --- that SSL-SGD creates Hessian sharpness anisotropy along spurious feature directions --- has not yet been empirically verified. The full 200-epoch run is ongoing. This is the most significant limitation: the paper motivates and designs a geometric intervention without yet demonstrating the geometric phenomenon it seeks to address. The WGA gap data independently establishes that a geometric problem exists; the measurement protocol is designed and validated at the code level. The theoretical motivation (Gatmiry 2024, SCER 2025, G2-SAM 2025) is well-grounded.

\textbf{Limitation 2: Scope limited to DINO + Waterbirds (confirmed results).}
The confirmed results come from a single SSL method (DINO) on a single dataset (Waterbirds). The planned evaluation covers 9 combinations (3 SSL methods $\times$ 3 datasets). CelebA was excluded due to a WILDS download error (HTTP 500); CMNIST and MoCo-v2/SimCLR were excluded from the fast run due to speed constraints. Waterbirds is the canonical spurious correlation benchmark, and the gap is large enough to be definitive for the existence question.

\textbf{Limitation 3: Linear probe proxy precision/recall unvalidated.}
The assumption that the top-25\% high-loss linear probe samples accurately identify spurious/minority group samples has not been validated via precision/recall against held-out Waterbirds group labels. LFR \citep{ghaznavi2023lfr} validated an equivalent proxy in supervised ERM settings; transfer to SSL representations requires explicit verification (required metric in Section 4.5). If the proxy fails in SSL representations, the anisotropy ratio measures nothing meaningful.

\textbf{Limitation 4: cuDNN disabled.}
Training was conducted with cuDNN disabled due to a torch+cu124 / CUDA 12.9 driver incompatibility. This slows training but does not affect numerical results.

\textbf{Limitation 5: SAM computational cost.}
SAM requires two forward-backward passes per batch ($2\times$ cost vs. SGD). For 200-epoch SSL pretraining on ResNet-50, this approximately doubles pretraining time.

\subsubsection{6.3 Theoretical Implications}

The central theoretical question this work raises --- but cannot yet resolve --- is whether Gatmiry et al.'s [2024] rank-1 simplicity bias result for supervised cross-entropy transfers to InfoNCE objectives. If it does, SAM would increase shortcut reliance in SSL (by promoting simpler/lower-rank features, which tend to be spurious). If the InfoNCE landscape prevents rank-1 dynamics (due to its uniform distribution pressure and multiple attractors), SAM's flattening effect would distribute more broadly and potentially reduce spurious anisotropy.

A finding that AR < 1.0 under SAM (SAM increases spurious anisotropy) would be a negative result of significant theoretical importance: it would demonstrate that the InfoNCE landscape responds to SAM differently than cross-entropy, ruling out optimizer geometry as the mechanism and pointing toward representational or objective-level explanations.

\subsubsection{6.4 Broader Impact}

SSL pretraining is increasingly deployed in fairness-sensitive domains: medical image analysis, hiring, content moderation. If SSL representations systematically encode spurious correlations geometrically, practitioners relying on these representations --- even with downstream fairness interventions --- face a structural problem. This work provides both a diagnostic (anisotropy measurement) and a potential intervention (SAM during pretraining) that could be applied before downstream deployment.

The annotation-free design is critical for practical adoption: acquiring group labels is expensive and requires domain expertise that may be unavailable in deployment settings.

